% =============================================================================
% DATOS DE LA TESIS
% =============================================================================

\newcommand{\thesistitle}{Diseño e Implementación de un Modelo de Evaluación para la Contención Selectiva en un Núcleo de Red 5G Cloud-Native}
\newcommand{\thesisauthor}{Denilson Ademir Quispe Ayala}
\newcommand{\thesisadvisor}{José Vicente de Paúl Rodríguez Salazar}

\title{\thesistitle}
\author{\thesisauthor}
\date{\limaFecha\today}

\hypersetup{
    pdftitle={\thesistitle},
    pdfauthor={\thesisauthor}
}

% =============================================================================
% COMANDOS DE ESTRUCTURA
% =============================================================================

% Capítulo no numerado que también aparece en el índice general.
\newcommand{\chapterwithtoc}[1]{%
    \chapter*{#1}
    \phantomsection
    \addcontentsline{toc}{chapter}{#1}%
}

% Nota de fuente uniforme para figuras y tablas.
\newcommand{\notafuente}[1]{%
    \caption*{Fuente: #1}%
}

% =============================================================================
% ESTILO DE LAS FIGURAS TÉCNICAS
% =============================================================================

% Paleta única para todos los diagramas de elaboración propia.
\definecolor{figNFborde}{HTML}{4A6D96}
\definecolor{figNFfondo}{HTML}{EDF2F8}
\definecolor{figNFtexto}{HTML}{1F3A5F}
\definecolor{figInfraBorde}{HTML}{7A7A7A}
\definecolor{figInfraFondo}{HTML}{F4F4F4}
\definecolor{figInfraTexto}{HTML}{333333}
\definecolor{figPEPborde}{HTML}{B58133}
\definecolor{figPEPfondo}{HTML}{FBF3E4}
\definecolor{figPEPtexto}{HTML}{6E4E17}
\definecolor{figNegBorde}{HTML}{B03A2E}
\definecolor{figNegFondo}{HTML}{FBEDEB}
\definecolor{figNegTexto}{HTML}{7B2A21}
\definecolor{figPosBorde}{HTML}{3B7A57}
\definecolor{figPosFondo}{HTML}{EDF5F0}
\definecolor{figPosTexto}{HTML}{2A5A40}
\definecolor{figGris}{HTML}{6E6E6E}
\definecolor{figGrisSuave}{HTML}{A8A8A8}
\definecolor{figBanda}{HTML}{FAFAFA}

\tikzset{
    % Base común: grosor de línea, punta de flecha y tipografía.
    figbase/.style={
        font=\small,
        line width=0.7pt,
        >={Latex[length=2.1mm,width=1.5mm]},
        every node/.append style={inner sep=3pt}
    },
    % Función de red.
    fignf/.style={
        draw=figNFborde, fill=figNFfondo, text=figNFtexto,
        rounded corners=2pt, line width=0.8pt,
        minimum width=2.4cm, minimum height=1cm, align=center
    },
    % Componente de infraestructura o intermediario.
    figinfra/.style={
        draw=figInfraBorde, fill=figInfraFondo, text=figInfraTexto,
        line width=0.9pt,
        minimum width=2.2cm, minimum height=1cm, align=center
    },
    % Punto de ejecución de políticas.
    figpep/.style={
        draw=figPEPborde, fill=figPEPfondo, text=figPEPtexto,
        shape=chamfered rectangle, chamfered rectangle xsep=7pt,
        line width=0.8pt,
        minimum width=2cm, minimum height=1cm, align=center
    },
    % Registro de instancias.
    fignrf/.style={
        draw=figInfraBorde, fill=figInfraFondo, text=figInfraTexto,
        shape=cylinder, shape border rotate=90, aspect=0.18,
        line width=0.8pt,
        minimum width=1cm, minimum height=1.9cm, align=center
    },
    % Comprobación de una cadena de decisión.
    figgate/.style={
        draw=figNFborde, fill=figNFfondo, text=figNFtexto,
        rounded corners=2pt, line width=0.8pt,
        minimum width=4.2cm, minimum height=1.05cm, align=center
    },
    % Resultado negativo y su causa.
    figneg/.style={
        draw=figNegBorde, fill=figNegFondo, text=figNegTexto,
        rounded corners=2pt, line width=0.8pt,
        align=left, text width=6.5cm, inner sep=6pt
    },
    % Veredicto positivo.
    figpos/.style={
        draw=figPosBorde, fill=figPosFondo, text=figPosTexto,
        rounded corners=2pt, line width=1pt,
        minimum width=4.2cm, minimum height=1.05cm, align=center
    },
    % Nota al margen del flujo principal.
    fignota/.style={
        draw=figGrisSuave, fill=white, text=figInfraTexto,
        line width=0.6pt, dash pattern=on 2pt off 1.6pt,
        align=left, inner sep=6pt, font=\footnotesize
    },
    % Aristas.
    figflujo/.style={->, draw=figGris, line width=0.9pt},
    figsi/.style={->, draw=figPosBorde, line width=1pt},
    figno/.style={->, draw=figNegBorde, line width=0.8pt,
                  dash pattern=on 3pt off 2pt},
    figaux/.style={->, draw=figGrisSuave, line width=0.6pt,
                   dash pattern=on 2pt off 1.6pt},
    % Rótulos.
    figrot/.style={font=\scriptsize, text=figGris, inner sep=2pt},
    figcarril/.style={font=\small\bfseries, text=figInfraTexto, anchor=west}
}
